\documentclass[conference]{IEEEtran}

\usepackage[utf8]{inputenc}
\usepackage[T1]{fontenc}
\usepackage{amsmath,amssymb}
\usepackage{graphicx}
\usepackage{booktabs}
\usepackage{multirow}
\usepackage[table]{xcolor}
\usepackage{hyperref}
\usepackage{balance}
\usepackage{subcaption}
\usepackage{siunitx}
\usepackage{microtype}
\usepackage[font=small,labelfont=bf,skip=4pt]{caption}
\usepackage{stfloats}
\usepackage{float}
\usepackage{enumitem}
\setlist{nosep,leftmargin=*}
\renewcommand{\arraystretch}{1.1}
\relpenalty=9999
\binoppenalty=9999

\hypersetup{
  colorlinks=true, linkcolor=black,
  citecolor=blue!60!black, urlcolor=blue!60!black,
}

\title{Ranking Inversion in Emergency UAV Planning:\\
       Why the Best Navigator Is Not the Best Rescuer}

\author{%
  \IEEEauthorblockN{Konstantinos~Zervakis\IEEEauthorrefmark{1}
    \and Elias~Panagiotopoulos\IEEEauthorrefmark{1}}
  \IEEEauthorblockA{\IEEEauthorrefmark{1}Athens, Greece}
}

\begin{document}
\maketitle

% ============================================================================
\begin{abstract}
Path-planning benchmarks rank UAV planners by navigation success rate---a
binary metric that ignores \emph{when} the goal is reached. We demonstrate
that this metric is not merely incomplete but \emph{actively misleading}:
in time-critical disaster response, the planner that survives most often
delivers the worst mission outcomes. We call this the \emph{ranking
inversion}. To establish this finding, we introduce FLARE, a deterministic
benchmark evaluating five risk-parameterised planners across
450~episodes on three OpenStreetMap-derived Greek wildfire scenarios with
coupled hazard dynamics (wind-driven fire, structural collapse, traffic
closures). Each planner differs by a single risk coefficient~$\rho$ that
controls sensitivity to a shared hazard cost map. FLARE measures mission
impact---medication efficacy, casualty survival, surveillance
freshness---all strictly decreasing functions of delivery time. The result:
the planner with the highest success rate (62.7\%) ranks \emph{last} among
adaptive planners in mission score ($\Delta$Rank$\,{=}\,{-3}$), while a
risk-tolerant planner achieves the highest mission score despite ranking
third in navigation ($\Delta$Rank$\,{=}\,{+2}$; Friedman $p < 0.001$).
The mechanism is precise: higher~$\rho$ induces paths with higher survival probability, longer
delivery time~$t$, and lower mission score~$M(t)$. Risk is not penalised directly---it is
penalised through time-induced mission degradation.
\end{abstract}

\begin{IEEEkeywords}
UAV path planning, disaster response, ranking inversion, risk-aware
planning, wildfire, mission-impact scoring, reproducibility
\end{IEEEkeywords}

% ============================================================================
\section{Introduction}\label{sec:intro}

% 1. REAL-WORLD HOOK
During the 2021 Evia wildfire---the largest recorded in the European Union
at the time---settlements were isolated for days and medical supplies had to
reach them by sea and air~\cite{Erdelj2017,Rahim2025}. Two years later, the
2023 Evros megafire consumed 96{,}000~hectares under complex no-fly-zone
constraints~\cite{copernicus_emsn166,effis2023}. In both events, UAVs could
have delivered medicine, surveyed fire perimeters, and located stranded
casualties---provided a planner could reach the target within a time horizon where
the delivery still has value.

% 2. PROBLEM GAP
Current benchmarks cannot evaluate this requirement. They reduce every
mission to a binary question---\emph{did the agent reach the goal?}---and
rank planners by success rate alone. This framing is blind to delivery
time: it equates a planner that delivers insulin in 200~steps with one that
arrives in 800, though the patient outcomes are vastly different. Navigation
success is not a proxy for mission effectiveness.

% 3. CORE QUESTION
This paper poses the question directly: \emph{does the planner that
survives most often also save the most lives?}

% 4. MAIN RESULT
Across 450~deterministic episodes on three Greek wildfire scenarios, the
answer is \textbf{no}. We term this phenomenon the \emph{ranking
inversion}: planner orderings by navigation success rate diverge from
orderings by mission impact score.

% 5. WHY IT HAPPENS
The cause traces to a single parameter---the risk
coefficient~$\rho \in \mathbb{R}_{\geq 0}$---that governs how strongly
each planner avoids hazards. Higher~$\rho$ amplifies perceived danger,
producing wider detours with higher survival probability.
Lower~$\rho$ discounts risk, permitting shorter-time routes through fire
fringes at the cost of increased collision probability. Risk is not penalised
directly---it is penalised through time-induced mission degradation:
\begin{equation*}
  \rho \;\uparrow\;\; \Rightarrow \;\;
  \text{higher } P(\text{survive}) \;\Rightarrow\;\;
  t \;\uparrow\;\; \Rightarrow \;\;
  M(t) \;\downarrow
\end{equation*}
where $t \in \mathbb{N}$ is the delivery step and $M: \mathbb{N} \to [0,1]$
is the strictly decreasing mission score.

% 6. CONTRIBUTIONS
\textbf{Contributions.}
\begin{itemize}
  \item \textbf{FLARE benchmark}: three OSM-derived wildfire scenarios with
    coupled hazard dynamics and time-decaying mission scores, grounded in
    real Greek incidents, enforcing 36~deterministic contracts.
  \item \textbf{Risk-parameterised planner suite}: five planners spanning
    static, adaptive, and reactive families, differing solely in~$\rho$.
  \item \textbf{Ranking inversion}: empirically demonstrated and
    statistically validated (Friedman $p < 0.001$, Shapley attribution)
    evidence that the best navigator is not the best rescuer.
\end{itemize}

The remainder is organised as follows: \S\ref{sec:related} positions FLARE
relative to existing work; \S\ref{sec:design} formalises the benchmark;
\S\ref{sec:results} presents experiments; \S\ref{sec:discussion} interprets
findings; \S\ref{sec:conclusion} concludes.

% ============================================================================
\section{Related Work}\label{sec:related}

\textbf{UAV path planning.} Classical planners---A*~\cite{Hart1968},
D*~Lite~\cite{Koenig2002}, APF~\cite{Khatib1986}---are surveyed
extensively~\cite{Meng2025,Dradoum2025,Aggarwal2020}. Risk-aware
extensions~\cite{Zhou2025,Melissaris2025} weight edges by hazard proximity
but evaluate only path length or success rate. No existing work evaluates
the \emph{mission-level} consequences of risk sensitivity---the gap that
FLARE addresses.

\textbf{Reproducible benchmarking.}
Henderson et al.~\cite{Henderson2018} demonstrated evaluation fragility
without fixed seeds. Robust practice now includes IQM with bootstrap
CIs~\cite{Agarwal2021}, Friedman/Wilcoxon tests~\cite{Demsar2006}, and
performance profiles~\cite{Dolan2002}. Bonsignorio et
al.~\cite{Bonsignorio2025} advocate published experimental contracts.
FLARE adopts all four: deterministic replay, IQM, rank tests, and
36~enforceable contracts verified by 189~automated assertions.

\textbf{Dynamic fire environments.} Cellular-automaton models with
wind-directional spread~\cite{Alexandridis2008,Giannino2025,Sahila2025}
capture the dominant fire propagation mechanism. Most UAV benchmarks,
however, use static or randomly toggled
blockages~\cite{Zhao2020}---environments that do not worsen over time.
FLARE's physically motivated dynamics (spreading fire, collapsing
buildings, closing roads) create a \emph{monotonically degrading}
environment, a qualitatively harder and more realistic challenge.

\textbf{The missing evaluation paradigm.} Across all three threads,
existing work shares one blind spot: \emph{time-coupled mission value}.
Success rate counts completions; FLARE measures their worth.

% ============================================================================
\section{Benchmark Design}\label{sec:design}

\subsection{Environment}

FLARE operates on a grid $\mathcal{G} = \{0,\dots,499\}^2$ (${\sim}3$\,m
per cell), extracted from OpenStreetMap tiles~\cite{Boeing2017} of Greek
urban areas. The Gymnasium-compatible~\cite{Towers2025} environment exposes
$|\mathcal{A}|{=}5$ actions (four cardinal moves and stay). A strict timing
contract governs each step: the planner observes the \emph{current}
blocking mask (frozen), selects an action, and only \emph{then} do dynamics
advance (fairness contract FC-2). If fire or debris reaches the agent's
cell after the dynamics update, the episode terminates immediately.

\subsection{Hazard Dynamics}\label{sec:dynamics}

Three coupled layers produce a hazard landscape that worsens monotonically.

\textbf{Wind-driven fire.} Spread follows a cellular automaton on the Moore
neighbourhood~\cite{Alexandridis2008}, modulated by wind direction and land
use. Downwind cells ignite faster; upwind cells are suppressed
(Fig.~\ref{fig:wind}). A buffer zone (radius~1) and dense smoke further restrict UAV movement.

\begin{figure*}[!t]
  \centering
  \caption{Wind reshapes the hazard geometry. (a)~Without wind, fire expands
    symmetrically. (b)~At 3\,m/s from the west, fire elongates downwind,
    blocking the direct route while the upwind flank remains passable for
    risk-tolerant planners. Black circles denote dynamic obstacles; blue
    cross denotes the UAV\@. Penteli, seed~42, $t{=}150$.}
  \label{fig:wind}
\end{figure*}

\textbf{Structural collapse.} Buildings exposed to fire for
$\tau_c{=}80$~steps collapse with probability $p_d{=}0.6$, scattering
permanent debris. Unlike fire---which burns out---debris
\emph{accumulates}, progressively eliminating detour options.
(Fig.~\ref{fig:collapse}).

\begin{figure}[t]
  \centering
  \caption{Cascading hazard coupling. (a)~Fire peaks then recedes, but
    debris grows permanently. (b--e)~Four snapshots trace the cascade
    from clean grid to permanent debris field (Penteli, seed~42).}
  \label{fig:collapse}
\end{figure}

\textbf{Dynamic obstacles.} Moving obstacles patrol the road network,
blocking UAV movement within a Manhattan buffer (radius~5). The
\texttt{InteractionEngine} couples fire to traffic: fire proximity closes
road segments, creating compound blockages that neither layer produces
alone (Fig.~\ref{fig:traffic}).

\begin{figure*}[!t]
  \centering
  \caption{Fire--obstacle coupling (Piraeus, seed~42). (a)~Dynamic
    obstacles operate normally at $t{=}30$. (b)~Fire triggers road
    closures by $t{=}80$. (c)~At $t{=}150$, fire, closures, and diverted
    obstacles jointly cut the corridor. Black circles denote dynamic
    obstacles; blue cross denotes the UAV\@.}
  \label{fig:traffic}
\end{figure*}

\subsection{Risk Cost Map}\label{sec:riskmap}

The perception pipeline maps raw hazards to planner-specific edge costs
in two stages: \emph{risk fusion} and \emph{cost inflation}.

\textbf{Stage 1: Risk fusion.} Five distance-decay layers are fused into a
continuous risk map $R: \mathcal{G} \to [0,1]$ via element-wise max:
\begin{equation}\label{eq:risk}
  R(\mathbf{x}) = \max\bigl(R_f,\; R_s,\; R_t,\; R_b,\; R_{\text{nfz}}\bigr)
\end{equation}
where $R_f{=}[1{-}d_f/r_f]^+$ (fire), $R_s{=}s/2$ (smoke),
$R_t{=}0.6[1{-}d_t/r_t]^+$ (traffic), $R_b{=}0.7[1{-}d_b/r_b]^+$
(debris), and $R_{\text{nfz}} \in \{0, 0.5, 1\}$ (no-fly zones).
for $\mathbf{x} \in \mathcal{G}$, where $d_f, d_t, d_b$ are Euclidean
distance transforms from fire, traffic, and debris with influence radii
$r_f{=}30$, $r_t{=}10$, $r_b{=}8$ cells; $s \in [0,1]$ is smoke
concentration; $[\cdot]^{+}{=}\max(\cdot,0)$. The scaling constants
(0.6, 0.7) are calibrated heuristics reflecting the relative severity
of each hazard layer (insensitive to $\pm 20\%$ perturbation). We use
$\max$ rather than summation so that a single severe hazard dominates
regardless of benign neighbours.

\emph{Intuition:} $R(\mathbf{x}){=}1$ means immediate danger;
$R(\mathbf{x}){=}0$ means no nearby hazard.

\textbf{Stage 2: Cost inflation.} Each planner inflates its edge cost
using a planner-specific risk coefficient $\rho \in \mathbb{R}_{\geq 0}$:
\begin{equation}\label{eq:cost}
  w(\mathbf{x}) = 1 + \rho \cdot R(\mathbf{x})
\end{equation}

\emph{Intuition:} High~$\rho$ amplifies perceived danger, causing wider
detours; low~$\rho$ allows routes through fire fringes.

\textbf{Summary.} The risk map encodes spatial danger; the coefficient
$\rho$ determines how strongly the planner reacts.
Fig.~\ref{fig:riskperception} establishes visually that $\rho$ alone---not
the algorithm---controls what the planner ``sees'': the \emph{same}
fire state produces five radically different cost landscapes.

\begin{figure*}[!t]
  \centering
  \caption{Same fire, five perceptions ($t{=}150$). Each panel shows
    $w(\mathbf{x}) = 1 + \rho \cdot R(\mathbf{x})$. A* sees uniform cost
    (blind). Periodic ($\alpha{=}5.0$) sees fire as an impassable wall.
    Aggressive ($\beta{=}0.5$) barely registers danger. The coefficient,
    not the algorithm, determines what the planner ``sees.''}
  \label{fig:riskperception}
\end{figure*}

\subsection{Planner Suite}\label{sec:planners}

Five planners from three algorithmic families are controlled by a single
coefficient~$\rho$ (Table~\ref{tab:planner_profiles}).
\textbf{A*}~\cite{Hart1968}: static baseline; plans once, $\rho{=}0$
(ignores risk).
\textbf{Periodic}: replans every 6~steps, $\alpha{=}5.0$ (risk-averse).
\textbf{Aggressive}: replans on mask change, $\beta{=}0.5$
(risk-tolerant).
\textbf{Incremental A*} (D*~Lite-inspired~\cite{Koenig2002}): replans on
path invalidation, $\gamma{=}2.0$.
\textbf{APF}~\cite{Khatib1986}: gradient repulsion, $\delta{=}3.0$.

\input{../outputs/paper_tables/planner_profiles.tex}

\subsection{Mission Scoring}\label{sec:missions}

Standard benchmarks count successes. FLARE measures the \emph{value} of
each success by defining a mission score $M(t) \in [0,1]$ that is
\emph{strictly decreasing} in delivery time~$t \in \mathbb{N}$
(Fig.~\ref{fig:missiondecay}). Three mission types specialise this
framework:

\textbf{Pharma delivery} (Penteli): medication efficacy decays
quadratically, $E(t) = \max(0,\, 1 - (t/T)^2)$, $T{=}T_{\max}$.

\emph{Intuition:} Insulin that arrives 200~steps early retains 94\%
efficacy; 400~steps late, only 75\%.

\textbf{Urban rescue} (Piraeus): casualty survival decays exponentially
with fire-coupled amplification:
\begin{equation}\label{eq:survival}
  S_i(t) = \exp\bigl(-\lambda_{\text{eff},i} \cdot t\bigr), \;\;
  \lambda_{\text{eff}} = \lambda_0
  \bigl(1 + \kappa / \max(d_{\text{fire}}, 1)\bigr)
\end{equation}
where $\lambda_0$ is the base decay rate and $\kappa{=}5.0$ is the
fire-coupling constant.

\emph{Intuition:} Casualties near active fire expire faster, explicitly
rewarding planners willing to enter dangerous zones.

\textbf{Fire surveillance} (Downtown): data freshness decays linearly,
$F(t) = 1 - t/T$.

\emph{Intuition:} A fire perimeter map becomes stale as the front advances.

\textbf{Key property.} All three functions share one characteristic: they
are strictly decreasing in~$t$. Every additional step of detour
\emph{destroys mission value}. Cross-scenario comparison uses
per-scenario normalisation (divide by scenario max, then average).

\begin{figure}[t]
  \centering
  \caption{Why speed matters. All three mission scores $M(t)$ decay with
    delivery time. The shaded gap between Aggressive (low~$t$) and Periodic
    (high~$t$) quantifies the $M(t)$ cost of a risk-averse detour.}
  \label{fig:missiondecay}
\end{figure}

\subsection{Risk--Time--Mission Coupling}\label{sec:coupling}

The link between risk aversion and mission failure is now explicit.
Increasing~$\rho$ produces safer paths with higher survival probability but
longer travel time~$t$. Because every mission score $M(t)$ is strictly
decreasing in~$t$, the additional travel time directly reduces mission
value. This creates a fundamental trade-off:
\[
  \underbrace{\rho \;\uparrow}_{\text{more risk-averse}}
  \;\;\Longrightarrow\;\;
  \underbrace{t \;\uparrow}_{\text{longer detour}}
  \;\;\Longrightarrow\;\;
  \underbrace{M(t) \;\downarrow}_{\text{worse mission outcome}}
\]
This coupling is the \emph{mechanism} behind the ranking inversion: the
planner with the highest survival probability is penalised not by
hazards but by the delivery time its detours impose---an inevitable
consequence of any strictly time-decreasing objective.

\subsection{Scenarios}\label{sec:scenarios}

Each scenario is grounded in a real Greek wildfire event
(Table~\ref{tab:scenario_overview}).
\textbf{Penteli} (Evia 2021): wind-driven fire in a wildland--urban
interface; the mission is insulin delivery before efficacy degrades.
\textbf{Piraeus} (Rhodes 2023): dense port with structural collapse;
three casualties with fire-coupled survival.
\textbf{Downtown} (Evros 2023): extreme building density (0.50) with
NFZ corridors for manned aircraft.

\input{../outputs/paper_tables/scenario_overview.tex}

\subsection{Reproducibility}

FLARE enforces 36~contracts across 15~families, verified by 189~automated
tests. Key guarantees: bit-identical replay from a single RNG seed (DC-1),
planner-agnostic corridor interdictions (FC-1), frozen hazard state during
planning (FD-3). Infeasible episodes are detected at reset and excluded
uniformly across all planners.

% ============================================================================
\section{Experiments and Results}\label{sec:results}

\subsection{Setup}

Five planners $\times$ three scenarios $\times$ 30~seeds yield
450~episodes; 7.8\% are structurally infeasible and excluded. Statistics
follow Agarwal et al.~\cite{Agarwal2021}: IQM with 95\% bootstrap CIs
(10{,}000~resamples), Wilcoxon signed-rank with Bonferroni correction,
Cliff's delta, and Friedman tests.

\subsection{Navigation Performance}

A* never succeeds: corridor dynamic obstacles block the reference path before fire
even arrives---a consequence of ignoring risk entirely.
Among adaptive planners, Incremental~A* leads (62.7\% feasible SR),
followed by Periodic (55.4\%), Aggressive (50.6\%), and APF (33.7\%).

Table~\ref{tab:mission_scorecard} disaggregates by scenario with
mission-specific metrics and verdicts. Two patterns stand out:
(i)~Piraeus yields the highest success rates (67--70\%) because its dense
layout provides many alternative routes;
(ii)~Downtown is hardest, with extreme density constraining detours.

\input{../outputs/paper_tables/mission_scorecard.tex}

\subsection{The Ranking Inversion}\label{sec:ranking}

Table~\ref{tab:planner_profiles} contains the paper's central result.
The $\Delta$~Rank column ($\text{Nav Rank} - \text{Mission Rank}$)
quantifies the discrepancy between navigation and rescue performance.
We define a \emph{ranking inversion} as the condition
$\arg\max_\pi \text{SR}(\pi) \neq \arg\max_\pi M(\pi)$,
where $\text{SR}$ is feasible success rate and $M$ is normalised mission
score. This condition holds across all three scenarios and is confirmed
by Friedman test ($p < 0.001$). Navigation success is \emph{not} a proxy
for mission effectiveness.

\textbf{The best navigator fails at rescue.} Incremental~A* ranks~\#1 in
SR but only~\#4 in mission score ($\Delta{=}{-3}$). Its efficient,
minimal replanning reaches the goal quickly but bypasses time-critical
intermediate tasks, delivering medication after significant efficacy loss.
\emph{Why:} low replan count (7.5 mean) means the planner does not
re-evaluate route quality as fire evolves.

\textbf{The best rescuer is not the safest.} Aggressive Replan ranks~\#3
in SR but~\#1 in mission score ($\Delta{=}{+2}$). Its low coefficient
($\beta{=}0.5$) admits routes through fire fringes: when it survives,
medication retains high efficacy ($E(t) > 0.7$) and casualties remain
alive ($S(t) > 0$). \emph{Why:} low~$\rho$ yields shorter paths
($t \downarrow$), directly increasing~$M(t)$.

The mechanism follows directly from the coupling in
\S\ref{sec:coupling}: Aggressive applies $w = 1 + 0.5R$, barely
inflating costs near fire, while Periodic applies $w = 1 + 5.0R$---a
$10\times$ stronger aversion signal that forces massive detours.
\emph{The detour saves the planner's life but costs the patient's.}
Fig.~\ref{fig:scatter}(a) makes this visible: planners on the
SR-optimal frontier (right side) fall \emph{below} the diagonal in mission
score---the two rankings disagree.

\begin{figure*}[!t]
  \centering
  \caption{The ranking inversion, visualised. (a)~SR vs.\ mission score;
    the dashed line connecting Incr.~A* to Aggressive crosses the
    diagonal---the two rankings disagree. (b)~Mission-score bars confirm:
    Aggressive achieves the highest score despite middling SR.}
  \label{fig:scatter}
\end{figure*}

\subsection{Qualitative Evidence}

Fig.~\ref{fig:comparison5} provides spatial confirmation. Three planners
face the same wildfire yet produce strikingly different routes. A* collides
with the first dynamic obstacle at $t{=}78$. Aggressive threads through smoke in
522~steps. Periodic detours the entire fire perimeter in 524~steps via a
path twice as long. The coefficient---not the algorithm---determines the route. Note that
Aggressive and Periodic arrive in nearly the same number of steps (522
vs.\ 524) but via radically different paths: $\rho$ controls route
geometry, not raw speed.

\begin{figure*}[!t]
  \centering
  \caption{Same fire, three fates (Penteli, seed~42). A* (left) collides
    with a dynamic obstacle. Aggressive (centre) threads through smoke.
    Periodic (right) detours safely---but insulin degrades en route. Black
    circles denote dynamic obstacles; blue cross denotes the UAV\@.
    Yellow cross: target. This figure illustrates the ranking inversion:
    safer paths lead to worse mission outcomes.}
  \label{fig:comparison5}
\end{figure*}

\subsection{Statistical Validation}

Wilcoxon tests (Bonferroni-corrected) confirm A* differs from every
adaptive planner ($p < 0.001$). Critically, \emph{no significant
difference} exists between Periodic and Aggressive in SR ($p = 1.0$):
navigation success alone cannot separate them. The discriminator is mission
score: a Friedman test yields $\chi^2(4) = 190.8$, $p < 0.001$,
establishing the ranking inversion as statistically robust
(Table~\ref{tab:stats_summary}).

\begin{table}[t]
\centering
\caption{Statistical validation. Friedman confirms ranking differences;
  bootstrap CIs quantify uncertainty.}
\label{tab:stats_summary}
\footnotesize
\begin{tabular}{@{}lrrl@{}}
\toprule
Test & Statistic & $p$ & Result \\
\midrule
Friedman (5 planners) & $\chi^2{=}190.8$ & ${<}0.001$ & Planners differ \\
A* vs.\ Adaptive (Wilcoxon) & --- & ${<}0.001$ & All significant \\
A* vs.\ Incr.\ A* (Cliff's $\delta$) & 0.58 & --- & Large effect \\
Incr.\ A* vs.\ APF (Cliff's $\delta$) & 0.27 & --- & Small effect \\
\midrule
\multicolumn{4}{@{}l}{\textit{IQM success rate, 95\% bootstrap CI
  (10{,}000 resamples):}} \\
Incr.\ A* & 62.7\% & \multicolumn{2}{l}{[50.9, 74.5]} \\
Periodic & 55.4\% & \multicolumn{2}{l}{[43.1, 67.8]} \\
Aggressive & 50.6\% & \multicolumn{2}{l}{[38.6, 62.6]} \\
APF & 33.7\% & \multicolumn{2}{l}{[22.6, 44.9]} \\
\bottomrule
\end{tabular}
\end{table}

\subsection{Hazard Ablation}

Isolating dynamics layers reveals their contributions. Fire is the
dominant differentiator: it reduces A* to 0\% while adaptive planners
retain 39--64\%. Traffic alone also eliminates A* but has near-zero
marginal impact on Periodic and Aggressive---their replanning fully
mitigates obstacle blockages. Combined dynamics reduce all planners by
10--20~pp relative to fire alone, confirming non-additive layer interaction.
Road closures grow proportionally with fire (via the
\texttt{InteractionEngine}); debris accumulates after the collapse delay
($\tau_c{=}80$). Aggressive Replan's timely detours maintain forward
progress; Periodic's risk-averse strategy produces extended stalls---a
direct consequence of the $\rho$--$t$--$M$ coupling.

% ============================================================================
\section{Discussion}\label{sec:discussion}

\subsection{Shapley Attribution}

Exact Shapley values~\cite{Shapley1953} over coalitions \{fire, traffic,
collapse\} confirm fire as the dominant hazard for all planners. Traffic
degrades A* ($\phi{=}-0.50$) and APF ($\phi{=}-0.25$) but has zero
marginal impact on Periodic and Aggressive---their replanning absorbs
obstacle blockages entirely. Collapse contributes near-zero marginal effect,
subsumed by the fire that triggers it.

\emph{Practical insight:} in fire-dominated scenarios, the risk coefficient
$\rho$ matters more than replanning frequency; in traffic-heavy
environments, any adaptive planner suffices.

\subsection{$\rho$ as a Deployment Knob}

The ranking inversion establishes that no universal ``best planner''
exists. Instead, $\rho$ provides a principled deployment knob:

\begin{itemize}
  \item \textbf{Rescue / medication}: $\rho \leq 0.5$. Minimise~$t$ to
    maximise~$M(t)$; accept higher collision probability.
  \item \textbf{Surveillance}: $\rho \geq 5.0$. Late data retains value;
    collision yields $M{=}0$.
\end{itemize}
This reframes planner selection from ``which algorithm wins?'' to ``what
risk level does this mission warrant?''---a question FLARE makes
empirically answerable.

\subsection{Limitations}

FLARE uses 2D grids without altitude or aerodynamics. The planners are
classical; the CA fire model is approximate. Three Athens scenarios provide
focused but not exhaustive coverage. Risk coefficients are fixed, not
adapted online. The inversion is not obvious \emph{a priori}: no prior benchmark includes
time-coupled scoring, and $|\Delta\text{Rank}|{=}3$ is not predictable
without experiments. Nor is it parameter tuning: $\rho$ is a structural
property of each planner's cost function. The inversion persists across all
three geographically distinct scenarios. These simplifications are deliberate:
the contribution is the \emph{evaluation paradigm} and the
\emph{ranking-inversion finding}, not the planners themselves.

% ============================================================================
\section{Conclusion}\label{sec:conclusion}

We presented FLARE, a benchmark that answers a question conventional
evaluations cannot pose: does the best navigator also rescue most
effectively? Across 450~episodes on three Greek wildfire scenarios, the
answer is no. The planner with the highest survival probability ranks last in
$M(t)$ because its detours increase~$t$ beyond the point where
delivery retains value. A risk-tolerant planner achieves the highest
mission score despite ranking third in navigation (Friedman $p < 0.001$).

This finding---the ranking inversion---demonstrates that navigation success
rate is not merely incomplete but actively misleading for time-critical
operations: it rewards the wrong planner. The implication is a paradigm shift: from
evaluating \emph{whether} planners succeed to evaluating \emph{when} they
deliver. The risk coefficient provides the deployment knob; FLARE provides
the evaluation platform. Both are open-source.

Code and data: \url{https://github.com/zervakisai/uavplan}.

% ============================================================================
\section*{Acknowledgments}

Map data \copyright~OpenStreetMap contributors (ODbL). Scenarios inspired
by the 2021 Evia wildfire~\cite{effis2021,copernicus_evia2021}, 2023
Rhodes fires~\cite{copernicus_emsr675}, and 2023 Evros
megafire~\cite{copernicus_emsn166,effis2023}. Data from
EFFIS~\cite{effis2021,effis2023} and Copernicus
EMS~\cite{copernicus_evia2021,copernicus_emsr675,copernicus_emsn166}.

\bibliographystyle{IEEEtran}
\bibliography{references}

\end{document}
